%%%%%%%%%%%%%%%%%%%%%%%%%%%%%%%%%%%%%%%%%%%%%%%%%%%%%%%%%%%%
% 14.17 GLOSSARY
% The Composition of Advanced Mathematics
%%%%%%%%%%%%%%%%%%%%%%%%%%%%%%%%%%%%%%%%%%%%%%%%%%%%%%%%%%%%

\section{Glossary}
\label{sec:glossary}

\prerequisite{
Familiarity with the notation and terminology introduced throughout
the textbook.
}

\learningobjectives{
By the end of this reference section, the reader should be able to:
\begin{itemize}
    \item locate common mathematical terms alphabetically;
    \item recall concise definitions of major concepts;
    \item distinguish closely related mathematical vocabulary;
    \item connect glossary terms with their canonical subject areas;
    \item use terminology consistently in mathematical writing.
\end{itemize}
}

%%%%%%%%%%%%%%%%%%%%%%%%%%%%%%%%%%%%%%%%%%%%%%%%%%%%%%%%%%%%
\subsection{Using the Glossary}
\label{subsec:using-glossary}
%%%%%%%%%%%%%%%%%%%%%%%%%%%%%%%%%%%%%%%%%%%%%%%%%%%%%%%%%%%%

This glossary provides concise definitions of recurring terms used throughout
the textbook.

The entries are intended as reference summaries rather than substitutes for
the complete definitions, theorems, examples, and proofs developed in the
main chapters.

%%%%%%%%%%%%%%%%%%%%%%%%%%%%%%%%%%%%%%%%%%%%%%%%%%%%%%%%%%%%
\subsection*{A}
%%%%%%%%%%%%%%%%%%%%%%%%%%%%%%%%%%%%%%%%%%%%%%%%%%%%%%%%%%%%

\begin{description}

\item[Abelian group.]
A group whose operation is commutative:

\[
ab=ba.
\]

\item[Absolute convergence.]
A series

\[
\sum a_n
\]

is absolutely convergent if

\[
\sum |a_n|
\]

converges.

\item[Absolute error.]
For exact value \(x\) and approximation \(\hat{x}\),

\[
\boxed{
|x-\hat{x}|.
}
\]

\item[Absolute value.]
For real \(x\),

\[
|x|
=
\begin{cases}
x, & x\geq0,\\
-x, & x<0.
\end{cases}
\]

\item[Acceleration.]
The derivative of velocity with respect to time:

\[
\mathbf{a}(t)
=
\mathbf{v}'(t)
=
\mathbf{r}''(t).
\]

\item[Adjacency matrix.]
A matrix representing which vertices of a graph are connected.

\item[Affine function.]
A function of form

\[
f(x)=Ax+b.
\]

\item[Algebraic number.]
A complex number satisfying a nonzero polynomial equation with rational
coefficients.

\item[Algorithm.]
A finite, well-defined procedure for solving a class of problems.

\item[Almost everywhere.]
True except on a set of measure zero.

\item[Alternating series.]
A series whose terms alternate in sign.

\item[Antiderivative.]
A function \(F\) satisfying

\[
F'=f.
\]

\item[Argument.]
For complex \(z\neq0\), the angle associated with its polar representation.

\item[Arithmetic mean.]
The ordinary average

\[
\bar{x}
=
\frac1n
\sum_{i=1}^{n}x_i.
\]

\item[Asymptote.]
A line or curve approached by another curve in a specified limiting sense.

\item[Asymptotic equivalence.]
The relation

\[
f(x)\sim g(x)
\]

meaning

\[
\frac{f(x)}{g(x)}
\to1
\]

in the stated limiting regime.

\item[Automorphism.]
An isomorphism from a mathematical structure to itself.

\end{description}

%%%%%%%%%%%%%%%%%%%%%%%%%%%%%%%%%%%%%%%%%%%%%%%%%%%%%%%%%%%%
\subsection*{B}
%%%%%%%%%%%%%%%%%%%%%%%%%%%%%%%%%%%%%%%%%%%%%%%%%%%%%%%%%%%%

\begin{description}

\item[Banach space.]
A complete normed vector space.

\item[Basis.]
A linearly independent spanning set of a vector space.

\item[Bayesian inference.]
Statistical inference in which prior information and observed data are
combined using probability models.

\item[Bijection.]
A function that is both injective and surjective.

\item[Binary operation.]
A rule combining two elements of a set to produce another element of that
set.

\item[Binomial coefficient.]

\[
\boxed{
\binom{n}{k}
=
\frac{n!}{k!(n-k)!}.
}
\]

\item[Boundary.]
The set of points arbitrarily close to both a set and its complement.

\item[Bounded function.]
A function whose values remain within a finite bound.

\item[Bounded sequence.]
A sequence \((a_n)\) for which there exists \(M>0\) such that

\[
|a_n|\leq M
\]

for all \(n\).

\item[Branch.]
A selected single-valued portion of a generally multivalued mathematical
relation, especially in complex analysis.

\end{description}

%%%%%%%%%%%%%%%%%%%%%%%%%%%%%%%%%%%%%%%%%%%%%%%%%%%%%%%%%%%%
\subsection*{C}
%%%%%%%%%%%%%%%%%%%%%%%%%%%%%%%%%%%%%%%%%%%%%%%%%%%%%%%%%%%%

\begin{description}

\item[Cardinality.]
The size of a set.

\item[Cauchy sequence.]
A sequence whose terms eventually become arbitrarily close to one another.

\item[Characteristic function.]
For random variable \(X\),

\[
\varphi_X(t)
=
\mathbb{E}[e^{itX}].
\]

\item[Characteristic polynomial.]
For square matrix \(A\),

\[
p_A(\lambda)
=
\det(\lambda I-A).
\]

\item[Closed set.]
A set whose complement is open.

\item[Closure.]
The smallest closed set containing a given set.

\item[Codomain.]
The declared target set of a function.

\item[Coefficient.]
A multiplicative factor attached to a mathematical term.

\item[Combination.]
An unordered selection from a collection.

\item[Compact.]
A topological space in which every open cover has a finite subcover.

\item[Complete.]
A metric space is complete if every Cauchy sequence converges to a point in
the space.

\item[Complex conjugate.]
For

\[
z=a+bi,
\]

the conjugate is

\[
\overline{z}
=
a-bi.
\]

\item[Complex number.]
A number of form

\[
a+bi,
\qquad
i^2=-1.
\]

\item[Composite number.]
An integer greater than \(1\) that is not prime.

\item[Composition.]
For functions \(f,g\),

\[
(f\circ g)(x)
=
f(g(x)).
\]

\item[Conditional probability.]

\[
\Pr(A\mid B)
=
\frac{
\Pr(A\cap B)
}{
\Pr(B)
},
\qquad
\Pr(B)>0.
\]

\item[Condition number.]
A measure of sensitivity of a mathematical problem to perturbations in input
data.

\item[Congruence.]
In number theory,

\[
a\equiv b\pmod n
\]

means \(n\mid(a-b)\).

In geometry, congruent objects have the same size and shape.

\item[Connected.]
A topological space that cannot be separated into two disjoint nonempty open
subsets.

\item[Constraint.]
A condition restricting the admissible solutions of a problem.

\item[Continuity.]
A function \(f\) is continuous at \(a\) if

\[
\lim_{x\to a}f(x)
=
f(a).
\]

\item[Contour.]
An oriented curve, commonly used in complex integration.

\item[Contraction.]
A map satisfying

\[
d(Tx,Ty)
\leq
q\,d(x,y),
\qquad
0\leq q<1.
\]

\item[Convergence.]
Approach toward a limiting mathematical object.

\item[Convex function.]
A function satisfying

\[
f((1-t)x+ty)
\leq
(1-t)f(x)+tf(y)
\]

for \(0\leq t\leq1\).

\item[Convex set.]
A set containing the line segment between every pair of its points.

\item[Correlation.]
A standardized measure of association between two variables.

\item[Covariance.]

\[
\operatorname{Cov}(X,Y)
=
\mathbb{E}
[
(X-\mu_X)(Y-\mu_Y)
].
\]

\item[Critical point.]
A point where a derivative or gradient vanishes or fails to exist, depending
on context.

\item[Curvature.]
A measure of how rapidly a curve or surface departs from being straight or
flat.

\end{description}

%%%%%%%%%%%%%%%%%%%%%%%%%%%%%%%%%%%%%%%%%%%%%%%%%%%%%%%%%%%%
\subsection*{D}
%%%%%%%%%%%%%%%%%%%%%%%%%%%%%%%%%%%%%%%%%%%%%%%%%%%%%%%%%%%%

\begin{description}

\item[Decision variable.]
A variable whose value is chosen in an optimization problem.

\item[Degree.]
A context-dependent measure including polynomial degree, vertex degree, field
extension degree, or differential-equation order-related concepts.

\item[Derivative.]
The local linear rate of change of a function.

\item[Determinant.]
A scalar associated with a square matrix that measures invertibility and
volume scaling.

\item[Diagonal matrix.]
A matrix whose off-diagonal entries are zero.

\item[Diagonalizable.]
A matrix \(A\) for which

\[
A=PDP^{-1}
\]

for some invertible \(P\) and diagonal \(D\).

\item[Differentiable.]
Possessing a derivative in the relevant sense.

\item[Differential equation.]
An equation involving an unknown function and one or more derivatives.

\item[Dimension.]
The number of vectors in a basis of a finite-dimensional vector space.

\item[Directed graph.]
A graph whose edges possess direction.

\item[Directional derivative.]
The rate of change of a scalar function in a specified direction:

\[
D_{\mathbf{u}}f
=
\nabla f\cdot\mathbf{u}
\]

for unit vector \(\mathbf{u}\).

\item[Discrete.]
Taking values from a finite or countable set, depending on context.

\item[Distribution.]
In probability, a mathematical description of how probability is assigned
to possible values of a random variable.

\item[Divergence.]
For vector field \(\mathbf{F}\),

\[
\nabla\cdot\mathbf{F}.
\]

The term also describes failure of a sequence or series to converge.

\item[Domain.]
The set of permitted inputs of a function.

\item[Dual problem.]
An optimization problem mathematically associated with a primal problem.

\item[Dynamic system.]
A mathematical model describing how a state changes over time.

\end{description}

%%%%%%%%%%%%%%%%%%%%%%%%%%%%%%%%%%%%%%%%%%%%%%%%%%%%%%%%%%%%
\subsection*{E}
%%%%%%%%%%%%%%%%%%%%%%%%%%%%%%%%%%%%%%%%%%%%%%%%%%%%%%%%%%%%

\begin{description}

\item[Eigenvalue.]
A scalar \(\lambda\) satisfying

\[
Av=\lambda v
\]

for some nonzero \(v\).

\item[Eigenvector.]
A nonzero vector \(v\) satisfying the eigenvalue equation.

\item[Element.]
An object belonging to a set.

\item[Entire function.]
A function holomorphic on all of \(\mathbb{C}\).

\item[Equation.]
A statement asserting equality between expressions.

\item[Equilibrium.]
A state that remains unchanged under the dynamics of a system.

\item[Equivalent.]
Related under a specified equivalence relation.

\item[Estimator.]
A rule or statistic used to estimate an unknown parameter.

\item[Estimate.]
A realized numerical value produced by an estimator.

\item[Euclidean norm.]
For \(x\in\mathbb{R}^n\),

\[
\|x\|_2
=
\sqrt{
\sum_i x_i^2
}.
\]

\item[Event.]
A subset of a probability sample space.

\item[Expected value.]
A probability-weighted average:

\[
\mathbb{E}[X].
\]

\item[Experimental design.]
The mathematical and statistical organization of how experimental data are
collected.

\item[Exponential growth.]
Growth proportional to the current amount, often modeled by

\[
y(t)=y_0e^{kt}.
\]

\end{description}

%%%%%%%%%%%%%%%%%%%%%%%%%%%%%%%%%%%%%%%%%%%%%%%%%%%%%%%%%%%%
\subsection*{F}
%%%%%%%%%%%%%%%%%%%%%%%%%%%%%%%%%%%%%%%%%%%%%%%%%%%%%%%%%%%%

\begin{description}

\item[Factor.]
A quantity multiplied by another quantity to produce a product.

\item[Factorial.]

\[
n!
=
n(n-1)\cdots1,
\qquad
0!=1.
\]

\item[Feasible point.]
A point satisfying every constraint in an optimization problem.

\item[Feasible region.]
The set of all feasible points.

\item[Field.]
A commutative ring in which every nonzero element has a multiplicative
inverse.

\item[Finite difference.]
A discrete approximation to a derivative formed from nearby function values.

\item[Fixed point.]
A point \(x\) satisfying

\[
f(x)=x.
\]

\item[Fourier series.]
A representation of a periodic function in sine, cosine, or complex
exponential modes.

\item[Fourier transform.]
An integral transform representing a function in frequency space.

\item[Function.]
A rule assigning every input in a domain exactly one output.

\end{description}

%%%%%%%%%%%%%%%%%%%%%%%%%%%%%%%%%%%%%%%%%%%%%%%%%%%%%%%%%%%%
\subsection*{G}
%%%%%%%%%%%%%%%%%%%%%%%%%%%%%%%%%%%%%%%%%%%%%%%%%%%%%%%%%%%%

\begin{description}

\item[Gamma function.]

\[
\Gamma(s)
=
\int_0^\infty
x^{s-1}e^{-x}\,dx.
\]

\item[Generating function.]
A power series encoding a sequence.

\item[Geodesic.]
A locally shortest or straightest path in a geometric space.

\item[Gradient.]
For scalar \(f\),

\[
\nabla f
=
\begin{bmatrix}
f_{x_1}\\
\vdots\\
f_{x_n}
\end{bmatrix}.
\]

\item[Graph.]
A collection of vertices together with edges connecting them.

\item[Greatest common divisor.]
The largest positive integer dividing two or more integers.

\item[Group.]
An algebraic structure with an associative operation, identity, and inverses.

\end{description}

%%%%%%%%%%%%%%%%%%%%%%%%%%%%%%%%%%%%%%%%%%%%%%%%%%%%%%%%%%%%
\subsection*{H}
%%%%%%%%%%%%%%%%%%%%%%%%%%%%%%%%%%%%%%%%%%%%%%%%%%%%%%%%%%%%

\begin{description}

\item[Hessian.]
The matrix of second partial derivatives of a scalar-valued function.

\item[Holomorphic.]
Complex differentiable on an open set.

\item[Homeomorphism.]
A continuous bijection with continuous inverse.

\item[Homomorphism.]
A map preserving specified algebraic operations.

\item[Hyperplane.]
An affine subspace of codimension one.

\item[Hypothesis.]
An assumption in a theorem or statistical test.

\end{description}

%%%%%%%%%%%%%%%%%%%%%%%%%%%%%%%%%%%%%%%%%%%%%%%%%%%%%%%%%%%%
\subsection*{I}
%%%%%%%%%%%%%%%%%%%%%%%%%%%%%%%%%%%%%%%%%%%%%%%%%%%%%%%%%%%%

\begin{description}

\item[Identity.]
An equation true for every value in its stated domain.

\item[Identity element.]
An element that leaves every element unchanged under the relevant operation.

\item[Image.]
The set of outputs attained by a function or map.

\item[Improper integral.]
An integral involving an infinite interval or unbounded integrand, interpreted
using limits.

\item[Independent events.]
Events \(A,B\) satisfying

\[
\Pr(A\cap B)
=
\Pr(A)\Pr(B).
\]

\item[Independent vectors.]
Vectors for which the zero linear combination has only trivial coefficients.

\item[Index.]
A subscript or other label used to identify members of an indexed family.

\item[Induction.]
A proof method involving a base case and inductive implication.

\item[Infimum.]
The greatest lower bound of a set.

\item[Injective.]
A function for which distinct inputs have distinct outputs.

\item[Inner product.]
A generalized dot product satisfying positivity, linearity, and symmetry or
conjugate symmetry.

\item[Integer.]
An element of

\[
\mathbb{Z}.
\]

\item[Integral.]
A mathematical representation of accumulation.

\item[Integral domain.]
A commutative ring with identity and no zero divisors.

\item[Interior.]
The largest open set contained in a given set.

\item[Inverse function.]
For bijective \(f\), the map \(f^{-1}\) satisfying

\[
f^{-1}(f(x))=x.
\]

\item[Invariant.]
A quantity or property unchanged under a specified transformation.

\item[Irrational number.]
A real number that cannot be represented as a ratio of integers.

\item[Isomorphism.]
A bijective structure-preserving map.

\end{description}

%%%%%%%%%%%%%%%%%%%%%%%%%%%%%%%%%%%%%%%%%%%%%%%%%%%%%%%%%%%%
\subsection*{J}
%%%%%%%%%%%%%%%%%%%%%%%%%%%%%%%%%%%%%%%%%%%%%%%%%%%%%%%%%%%%

\begin{description}

\item[Jacobian.]
The matrix of first partial derivatives of a multivariable map.

\item[Jacobian determinant.]
The determinant of a square Jacobian matrix, often describing local volume
scaling under a coordinate transformation.

\end{description}

%%%%%%%%%%%%%%%%%%%%%%%%%%%%%%%%%%%%%%%%%%%%%%%%%%%%%%%%%%%%
\subsection*{K}
%%%%%%%%%%%%%%%%%%%%%%%%%%%%%%%%%%%%%%%%%%%%%%%%%%%%%%%%%%%%

\begin{description}

\item[Kernel.]
For a linear or algebraic map, the set of elements mapped to zero or the
identity.

\item[Kronecker delta.]

\[
\delta_{ij}
=
\begin{cases}
1, & i=j,\\
0, & i\neq j.
\end{cases}
\]

\end{description}

%%%%%%%%%%%%%%%%%%%%%%%%%%%%%%%%%%%%%%%%%%%%%%%%%%%%%%%%%%%%
\subsection*{L}
%%%%%%%%%%%%%%%%%%%%%%%%%%%%%%%%%%%%%%%%%%%%%%%%%%%%%%%%%%%%

\begin{description}

\item[Lagrangian.]
A function used in mechanics, optimization, and variational mathematics,
with meaning determined by context.

\item[Laplacian.]
For scalar function \(f\),

\[
\Delta f
=
\nabla^2f.
\]

In graph theory, the graph Laplacian is commonly

\[
L=D-A.
\]

\item[Laplace transform.]

\[
\mathcal{L}\{f(t)\}
=
\int_0^\infty
e^{-st}f(t)\,dt.
\]

\item[Least squares.]
A method minimizing a sum of squared residuals.

\item[Limit.]
The value approached by a function, sequence, or other object in a specified
limiting process.

\item[Linear combination.]

\[
c_1v_1+\cdots+c_nv_n.
\]

\item[Linear independence.]
A set of vectors is linearly independent when only the trivial linear
combination equals zero.

\item[Linear map.]
A map satisfying

\[
T(au+bv)
=
aT(u)+bT(v).
\]

\item[Logarithm.]
The inverse operation of exponentiation.

\item[Local maximum.]
A point where the function is no smaller than nearby values.

\item[Local minimum.]
A point where the function is no larger than nearby values.

\end{description}

%%%%%%%%%%%%%%%%%%%%%%%%%%%%%%%%%%%%%%%%%%%%%%%%%%%%%%%%%%%%
\subsection*{M}
%%%%%%%%%%%%%%%%%%%%%%%%%%%%%%%%%%%%%%%%%%%%%%%%%%%%%%%%%%%%

\begin{description}

\item[Manifold.]
A space locally resembling Euclidean space.

\item[Matrix.]
A rectangular array of mathematical entries.

\item[Maximum.]
The greatest attained value of a set or function.

\item[Mean.]
A measure of center, commonly the arithmetic average.

\item[Measure.]
A generalized notion of size assigned to sets.

\item[Median.]
A central quantile dividing ordered data into two halves.

\item[Meromorphic.]
Holomorphic except at isolated poles.

\item[Metric.]
A distance function satisfying positivity, symmetry, and the triangle
inequality.

\item[Metric space.]
A set equipped with a metric.

\item[Minimum.]
The least attained value of a set or function.

\item[Model.]
A mathematical representation of selected features of a system.

\item[Modular arithmetic.]
Arithmetic with equivalence classes under congruence modulo an integer.

\item[Moment.]
A quantity such as

\[
\mathbb{E}[X^k].
\]

\item[Monotone.]
Consistently nondecreasing or nonincreasing.

\item[Monte Carlo method.]
A computational method based on random sampling.

\end{description}

%%%%%%%%%%%%%%%%%%%%%%%%%%%%%%%%%%%%%%%%%%%%%%%%%%%%%%%%%%%%
\subsection*{N}
%%%%%%%%%%%%%%%%%%%%%%%%%%%%%%%%%%%%%%%%%%%%%%%%%%%%%%%%%%%%

\begin{description}

\item[Natural number.]
A counting number, with inclusion of \(0\) depending on convention.

\item[Neighborhood.]
A set containing an open region around a point.

\item[Norm.]
A function measuring vector magnitude and satisfying the norm axioms.

\item[Normal distribution.]
The probability distribution

\[
N(\mu,\sigma^2)
\]

with bell-shaped density.

\item[Normal matrix.]
A complex matrix satisfying

\[
AA^*
=
A^*A.
\]

\item[Normal subgroup.]
A subgroup invariant under conjugation.

\item[Null hypothesis.]
The reference hypothesis in a statistical significance test.

\item[Null space.]
For matrix \(A\),

\[
\ker A
=
\{x:Ax=0\}.
\]

\item[Numerical method.]
An algorithm for approximating mathematical quantities computationally.

\end{description}

%%%%%%%%%%%%%%%%%%%%%%%%%%%%%%%%%%%%%%%%%%%%%%%%%%%%%%%%%%%%
\subsection*{O}
%%%%%%%%%%%%%%%%%%%%%%%%%%%%%%%%%%%%%%%%%%%%%%%%%%%%%%%%%%%%

\begin{description}

\item[Objective function.]
A function to be minimized or maximized.

\item[Open set.]
A member of the topology of a topological space.

\item[Operator.]
A map acting on mathematical objects, often functions or vectors.

\item[Optimal solution.]
A feasible point achieving the best objective value.

\item[Order.]
A context-dependent concept including relational order, group size,
differential-equation order, or convergence order.

\item[Ordinary differential equation.]
A differential equation involving derivatives with respect to one independent
variable.

\item[Orthogonal.]
Having zero inner product.

\item[Orthonormal.]
Orthogonal with unit norm.

\item[Outcome.]
An elementary result of a probability experiment.

\end{description}

%%%%%%%%%%%%%%%%%%%%%%%%%%%%%%%%%%%%%%%%%%%%%%%%%%%%%%%%%%%%
\subsection*{P}
%%%%%%%%%%%%%%%%%%%%%%%%%%%%%%%%%%%%%%%%%%%%%%%%%%%%%%%%%%%%

\begin{description}

\item[Parameter.]
A quantity characterizing a mathematical or statistical model and treated as
fixed within a specified analysis.

\item[Partial derivative.]
A derivative with respect to one variable while others are held fixed.

\item[Partial differential equation.]
A differential equation involving partial derivatives.

\item[Path.]
A continuous curve or a sequence of connected graph vertices, depending on
context.

\item[Permutation.]
An ordered arrangement.

\item[Polynomial.]
An expression formed from finitely many nonnegative integer powers with
coefficients.

\item[Power series.]
A series

\[
\sum_{n=0}^{\infty}
c_n(x-a)^n.
\]

\item[Prime.]
A positive integer greater than \(1\) with exactly two positive divisors.

\item[Probability.]
A measure of uncertainty satisfying the probability axioms.

\item[Probability density.]
A nonnegative function whose integral over a region gives continuous
probability.

\item[Probability mass.]
A probability assigned directly to a discrete outcome or value.

\item[Proof.]
A logically valid argument establishing a mathematical statement.

\item[Proposition.]
A mathematical statement presented for proof.

\item[\(p\)-value.]
A null-model tail probability measuring incompatibility between observed data
and the specified null hypothesis.

\end{description}

%%%%%%%%%%%%%%%%%%%%%%%%%%%%%%%%%%%%%%%%%%%%%%%%%%%%%%%%%%%%
\subsection*{Q}
%%%%%%%%%%%%%%%%%%%%%%%%%%%%%%%%%%%%%%%%%%%%%%%%%%%%%%%%%%%%

\begin{description}

\item[Quadratic form.]
An expression

\[
x^TAx.
\]

\item[Quantile.]
A value dividing a probability distribution or ordered data according to a
specified cumulative proportion.

\item[Quotient.]
A result of division or a mathematical object formed by identifying elements
under an equivalence relation.

\end{description}

%%%%%%%%%%%%%%%%%%%%%%%%%%%%%%%%%%%%%%%%%%%%%%%%%%%%%%%%%%%%
\subsection*{R}
%%%%%%%%%%%%%%%%%%%%%%%%%%%%%%%%%%%%%%%%%%%%%%%%%%%%%%%%%%%%

\begin{description}

\item[Random variable.]
A measurable numerical function defined on a probability space.

\item[Range.]
The set of outputs actually attained by a function.

\item[Rank.]
For a matrix or linear map, the dimension of its image.

\item[Rate of change.]
The amount one quantity changes relative to another.

\item[Rational number.]
A number expressible as

\[
\frac pq
\]

with integers \(p,q\) and \(q\neq0\).

\item[Real number.]
An element of

\[
\mathbb{R}.
\]

\item[Recurrence relation.]
An equation defining sequence terms using preceding terms.

\item[Regression.]
A statistical method relating a response variable to one or more explanatory
variables.

\item[Relative error.]

\[
\frac{
|x-\hat{x}|
}{
|x|
}
\]

for \(x\neq0\).

\item[Residual.]
Observed value minus fitted or predicted value.

\item[Residue.]
The coefficient of

\[
(z-z_0)^{-1}
\]

in a Laurent expansion.

\item[Ring.]
An algebraic structure with addition and multiplication satisfying the ring
axioms.

\item[Root.]
A value \(r\) satisfying

\[
f(r)=0.
\]

\end{description}

%%%%%%%%%%%%%%%%%%%%%%%%%%%%%%%%%%%%%%%%%%%%%%%%%%%%%%%%%%%%
\subsection*{S}
%%%%%%%%%%%%%%%%%%%%%%%%%%%%%%%%%%%%%%%%%%%%%%%%%%%%%%%%%%%%

\begin{description}

\item[Sample.]
Observed data or a collection drawn from a population or model.

\item[Sample space.]
The set of all possible outcomes of a probability experiment.

\item[Scalar.]
An element of the field underlying a vector space.

\item[Sequence.]
An ordered indexed collection

\[
(a_n).
\]

\item[Series.]
A sum of sequence terms.

\item[Set.]
A collection of distinct mathematical objects.

\item[Simulation.]
Computational generation of system behavior from a mathematical model.

\item[Singular matrix.]
A square matrix that is not invertible.

\item[Singularity.]
A point where a mathematical object fails to satisfy the regular behavior
under consideration.

\item[Span.]
The collection of all linear combinations of a set of vectors.

\item[Spectral radius.]
For square matrix \(A\),

\[
\rho(A)
=
\max_i|\lambda_i|.
\]

\item[Standard deviation.]
The square root of variance.

\item[Standard error.]
The standard deviation of an estimator or its estimate.

\item[State variable.]
A variable describing the current state of a dynamical system.

\item[Statistic.]
A numerical function of observed sample data.

\item[Stochastic process.]
A collection of random variables indexed by time, space, or another
parameter.

\item[Subgroup.]
A subset of a group that is itself a group under the inherited operation.

\item[Subset.]
A set \(A\) satisfying

\[
A\subseteq B.
\]

\item[Supremum.]
The least upper bound.

\item[Surjective.]
A function whose range equals its codomain.

\item[Symmetric matrix.]
A real matrix satisfying

\[
A^T=A.
\]

\end{description}

%%%%%%%%%%%%%%%%%%%%%%%%%%%%%%%%%%%%%%%%%%%%%%%%%%%%%%%%%%%%
\subsection*{T}
%%%%%%%%%%%%%%%%%%%%%%%%%%%%%%%%%%%%%%%%%%%%%%%%%%%%%%%%%%%%

\begin{description}

\item[Tangent.]
A line, vector, or space describing first-order local direction.

\item[Taylor series.]

\[
f(x)
=
\sum_{n=0}^{\infty}
\frac{
f^{(n)}(a)
}{
n!
}
(x-a)^n
\]

when the expansion converges to \(f\).

\item[Tensor.]
A multilinear mathematical object generalizing scalars, vectors, and
matrices.

\item[Theorem.]
A mathematical statement established by proof.

\item[Topology.]
A collection of open sets satisfying the topology axioms, or the subject
studying properties preserved by continuous deformation.

\item[Trace.]
For square matrix \(A\),

\[
\operatorname{tr}(A)
=
\sum_i a_{ii}.
\]

\item[Transform.]
A map converting a mathematical object into another representation.

\item[Transcendental number.]
A number that is not algebraic.

\item[Tree.]
A connected graph with no cycles.

\item[Triangle inequality.]

\[
\|x+y\|
\leq
\|x\|+\|y\|.
\]

\end{description}

%%%%%%%%%%%%%%%%%%%%%%%%%%%%%%%%%%%%%%%%%%%%%%%%%%%%%%%%%%%%
\subsection*{U}
%%%%%%%%%%%%%%%%%%%%%%%%%%%%%%%%%%%%%%%%%%%%%%%%%%%%%%%%%%%%

\begin{description}

\item[Uniform convergence.]
A sequence of functions \(f_n\) converges uniformly to \(f\) if

\[
\sup_x
|f_n(x)-f(x)|
\to0.
\]

\item[Unit vector.]
A vector with norm \(1\).

\item[Universal set.]
An ambient set relative to which complements or subsets are being discussed.

\end{description}

%%%%%%%%%%%%%%%%%%%%%%%%%%%%%%%%%%%%%%%%%%%%%%%%%%%%%%%%%%%%
\subsection*{V}
%%%%%%%%%%%%%%%%%%%%%%%%%%%%%%%%%%%%%%%%%%%%%%%%%%%%%%%%%%%%

\begin{description}

\item[Validation.]
Assessment of whether a model adequately represents the intended system for
its purpose.

\item[Variable.]
A symbol representing a quantity that may vary.

\item[Variance.]

\[
\operatorname{Var}(X)
=
\mathbb{E}
[
(X-\mathbb{E}[X])^2
].
\]

\item[Vector.]
An element of a vector space.

\item[Vector field.]
A function assigning a vector to each point of a domain.

\item[Verification.]
Assessment of whether a mathematical or computational model has been
implemented or solved correctly.

\item[Vertex.]
A node of a graph.

\end{description}

%%%%%%%%%%%%%%%%%%%%%%%%%%%%%%%%%%%%%%%%%%%%%%%%%%%%%%%%%%%%
\subsection*{W}
%%%%%%%%%%%%%%%%%%%%%%%%%%%%%%%%%%%%%%%%%%%%%%%%%%%%%%%%%%%%

\begin{description}

\item[Weighted graph.]
A graph whose edges or vertices are assigned numerical weights.

\item[Weighted mean.]

\[
\bar{x}_w
=
\frac{
\sum_i w_ix_i
}{
\sum_iw_i
}.
\]

\item[Work.]
In mechanics, the line integral of force along displacement:

\[
W
=
\int_C
\mathbf{F}\cdot d\mathbf{r}.
\]

\end{description}

%%%%%%%%%%%%%%%%%%%%%%%%%%%%%%%%%%%%%%%%%%%%%%%%%%%%%%%%%%%%
\subsection*{X}
%%%%%%%%%%%%%%%%%%%%%%%%%%%%%%%%%%%%%%%%%%%%%%%%%%%%%%%%%%%%

\begin{description}

\item[\(x\)-intercept.]
A point where a graph intersects the \(x\)-axis, satisfying

\[
y=0.
\]

\end{description}

%%%%%%%%%%%%%%%%%%%%%%%%%%%%%%%%%%%%%%%%%%%%%%%%%%%%%%%%%%%%
\subsection*{Y}
%%%%%%%%%%%%%%%%%%%%%%%%%%%%%%%%%%%%%%%%%%%%%%%%%%%%%%%%%%%%

\begin{description}

\item[\(y\)-intercept.]
A point where a graph intersects the \(y\)-axis, commonly obtained by setting

\[
x=0.
\]

\end{description}

%%%%%%%%%%%%%%%%%%%%%%%%%%%%%%%%%%%%%%%%%%%%%%%%%%%%%%%%%%%%
\subsection*{Z}
%%%%%%%%%%%%%%%%%%%%%%%%%%%%%%%%%%%%%%%%%%%%%%%%%%%%%%%%%%%%

\begin{description}

\item[Zero.]
The additive identity.

\item[Zero divisor.]
A nonzero element \(a\) of a ring for which some nonzero \(b\) satisfies

\[
ab=0.
\]

\item[Zero of a function.]
A point \(x\) satisfying

\[
f(x)=0.
\]

\item[\(z\)-score.]
A standardized value

\[
z
=
\frac{
x-\mu
}{
\sigma
}.
\]

\item[\(z\)-transform.]
For sequence \(x[n]\),

\[
X(z)
=
\sum_n
x[n]z^{-n}.
\]

\end{description}

%%%%%%%%%%%%%%%%%%%%%%%%%%%%%%%%%%%%%%%%%%%%%%%%%%%%%%%%%%%%
\subsection{Glossary Usage Notes}
\label{subsec:glossary-usage-notes}
%%%%%%%%%%%%%%%%%%%%%%%%%%%%%%%%%%%%%%%%%%%%%%%%%%%%%%%%%%%%

Some mathematical words have different meanings in different subjects.

Examples include:

\[
\boxed{
\text{normal},
\quad
\text{degree},
\quad
\text{order},
\quad
\text{rank},
\quad
\text{field},
\quad
\text{ring}.
}
\]

The intended meaning should always be determined from context.

%%%%%%%%%%%%%%%%%%%%%%%%%%%%%%%%%%%%%%%%%%%%%%%%%%%%%%%%%%%%
\subsection{Section Connections}
\label{subsec:glossary-connections}
%%%%%%%%%%%%%%%%%%%%%%%%%%%%%%%%%%%%%%%%%%%%%%%%%%%%%%%%%%%%

\chapterconnection{
The Glossary provides an alphabetical reference for terminology used across
the entire textbook. Mathematical Terminology in Section~14.16 groups related
terms by subject, while this section emphasizes rapid alphabetical lookup.
The final section, Master Index, organizes important concepts by topic and
identifies where they occur throughout the mathematical hierarchy.
}

%%%%%%%%%%%%%%%%%%%%%%%%%%%%%%%%%%%%%%%%%%%%%%%%%%%%%%%%%%%%
% SECTION SOURCE TRACKING
%%%%%%%%%%%%%%%%%%%%%%%%%%%%%%%%%%%%%%%%%%%%%%%%%%%%%%%%%%%%

%------------------------------------------------------------
% 14.17 Glossary
%------------------------------------------------------------
%
% Source basis:
%   - Standard terminology used throughout undergraduate
%     and graduate mathematics.
%
% Used for:
%   - alphabetical terminology reference;
%   - cross-disciplinary mathematical vocabulary;
%   - notation interpretation;
%   - concise definition lookup.
%
% Notes:
%   - Complete definitions remain in their canonical chapters.
%   - Master Index follows in Section 14.18.
%
%%%%%%%%%%%%%%%%%%%%%%%%%%%%%%%%%%%%%%%%%%%%%%%%%%%%%%%%%%%%